\begin{definition}
For integers $n\ge0$ and $k\ge0$ and a real number $p$ put
\[
  \epsilon_k(n,p)=\frac{\lambda_k(n,p)+n-k}{\lambda_k(n,p)+n+k},
\]
where $\lambda_k(n,p)$ is the exponent of \noderef{LambdaCoef}.  This is
equation (21) of the paper.

The quotient is interpreted as $0$ whenever its denominator vanishes.  The
paper uses $\epsilon_k(n,p)$ only in the range $1<p<\infty$, $k\ge n\ge1$,
where $\lambda_k(n,p)\ge0$ and hence $\lambda_k(n,p)+n+k\ge 2n>0$, so that
convention is never invoked.
\end{definition}
